\chapter{Harmonic candidate $\beta_N = 12(H_N - 1)$: scaling-law frontier
for every $N \ge 2$}
\label{ch:beta-N-closed-form-all}

\section*{Statement of the frontier}

The finite-$\beta_N$ criterion for $\cW_N$ tempering
(Chapter~\ref{ch:wn-tempered-closure},
Theorem~\ref{thm:wn-tempered-all-N}) is proved.
The low-rank data $\beta_2 = 6$ and $\beta_3 = 10$ admit two
polynomial extrapolations: the triangular formula
$\beta_N = (N+1)(N+2)/2$ and the quadratic formula
$\beta_N = N^2 - N + 4$. A single scaling-law conjecture distinguishes
these from the harmonic candidate at every $N \ge 4$.

\emph{Harmonic closed-form candidate
(Conjecture~\ref{conj:beta-N-harmonic-closed-form}).} For every
integer $N \ge 2$,
\[
 \beta_N \;=\; 12 \cdot (H_N - 1)
 \;=\; 12 \sum_{j = 2}^{N} \frac{1}{j},
\]
where $H_N = \sum_{j = 1}^{N} 1/j$ is the $N$-th harmonic number.
The values
\[
 \beta_2 = 6, \quad \beta_3 = 10, \quad
 \beta_4 = 13, \quad \beta_5 = \tfrac{77}{5}, \quad
 \beta_6 = \tfrac{87}{5}, \quad \ldots
\]
separate from the triangular and quadratic polynomials at $N = 4$
($15$ and $16$ versus $13$); under the harmonic law $\beta_2,\beta_3,
\beta_4$ are integral and $\beta_N$ is non-integral for every
$N\geq 5$. On the leading-root sharp locus, the conditional convergence
radius is
\[
 \rho_*^{\cW_N}(c)
 \;=\;
 \frac{\lvert c \rvert}{\beta_N}
 \;=\;
 \frac{\lvert c \rvert}{12(H_N - 1)}
 \;=\;
 \frac{\lvert c \rvert}{12 \sum_{j = 2}^{N} 1/j}.
\]
Theorem~\ref{thm:wn-tempered-all-N} uses only finiteness of $\beta_N$
and Stirling dominance after the Riccati envelope is supplied; the
harmonic closed form gives the conditional exact convergence radius only
after leading-root sharpness is imposed.

The supporting data are the Vol~II datum $A_4^{\cW_3} = 10/3$
(Proposition~\ref{prop:w3-shadow-leading-asymptotic}(3)) and the Vol~I
Virasoro closed form $A_r(\mathrm{Vir}) = 8(-6)^{r-4}/r$
(\textsf{thm:shadow-tower-asymptotic-closed-form} of
\texttt{shadow\_tower\_higher\_coefficients.tex}). Both are consistent
with the scaling law
$A_r(\cW_N) = (\kappaChHodge(\cW_N)/\kappaChHodge(\mathrm{Vir}))^{r-3} \cdot
A_r(\mathrm{Vir})$. The all-rank proof obligation is the direct
computation of $A_5^{\cW_4}/A_4^{\cW_4}$ from the full $\cW_4$ Miura
OPE.

\section{The $\kappa$-ratio scaling law}
\label{sec:kappa-ratio-scaling-law}

The load-bearing structural conjecture is that the full-shadow leading-Laurent
coefficient at the $\cW_N$ level scales from the Virasoro level by the
$\kappa$-ratio raised to the power $r-3$. Vol II's
Proposition~\ref{prop:w3-shadow-leading-asymptotic}(3) datum
$A_4^{\cW_3} = 10/3$ is its first non-Virasoro consistency check. If
the conjecture holds, it determines $\beta_N$ at all $N$.

\begin{conjecture}[$\kappa$-ratio scaling law]
\label{conj:kappa-ratio-scaling-law}
\ClaimStatusConjectured
For every integer $N \ge 2$ and every $r \ge 3$ with $r \ne 4$ mod parity
channels vanishing, the leading-Laurent coefficient of the full
$\cW_N$ shadow tower satisfies
\begin{equation}
\label{eq:kappa-ratio-scaling}
A_r(\cW_N)
\;=\;
\Bigl( \frac{\kappaChHodge(\cW_N)}{\kappaChHodge(\mathrm{Vir})} \Bigr)^{r-3}
\cdot A_r(\mathrm{Vir}),
\end{equation}
where $A_r(\mathrm{Vir}) = 8(-6)^{r-4}/r$ is the Virasoro value
(Vol~I \texttt{shadow\_tower\_higher\_coefficients.tex}
thm:shadow-tower-asymptotic-closed-form) and
$\kappaChHodge(\cW_N)/\kappaChHodge(\mathrm{Vir}) = 2(H_N - 1)$.
\end{conjecture}

\begin{remark}[Evidence and proof obligation for the scaling law]
Two facts are proved; one bridge remains open.

\emph{First proved datum.} By
Proposition~\ref{prop:w3-shadow-leading-asymptotic}(3) of Vol~II,
$A_4^{\cW_3} = 10/3$. The Virasoro value is
$A_4(\mathrm{Vir}) = 2$, hence
\[
 \frac{A_4^{\cW_3}}{A_4(\mathrm{Vir})}
 =
 \frac{10/3}{2}
 =
 \frac{5}{3}
 =
 \frac{\kappaChHodge(\cW_3)}{\kappaChHodge(\mathrm{Vir})}.
\]
This verifies \eqref{eq:kappa-ratio-scaling} at the single
non-Virasoro point $(N,r)=(3,4)$.

\emph{Second proved datum.} The direct finite-rank $\cW_4$ lane
calculation uses the top OPE norms \(c/2,c/3,c/4\) for the spin
\(2,3,4\) lanes. Its lane contribution is
\[
 6+4+3=13,
\]
separating the harmonic value from the triangular value \(15\) and the
quadratic value \(16\). This is a spin-lane witness, not a Riccati
bridge.

The full shadow recursion still has the schematic master equation
\begin{equation}
\label{eq:master-eqn-class-m}
S_r(\cW_N)
\;=\;
-\frac{1}{2 r \cdot \kappaChHodge(\cW_N)}
\sum_{\substack{j + k = r + 2 \\ 3 \le j \le k < r}}
 f(j, k) \cdot j \cdot k \cdot S_j(\cW_N) \cdot S_k(\cW_N).
\end{equation}
To turn the \(\cW_4\) lane witness into a proof of
\(\beta_4=13\), one must compute the leading coefficients
\(A_4^{\cW_4}\) and \(A_5^{\cW_4}\) from the full Miura/OPE constants in
\eqref{eq:master-eqn-class-m}, including the cross-channel and composite
terms, and verify
\begin{equation}
\label{eq:w4-missing-riccati-bridge}
 \frac{A_5^{\cW_4}}{A_4^{\cW_4}}
 =
 -\frac{52}{5}.
\end{equation}
Equivalently, if one uses the \(\kappa\)-scaling prediction
\(A_4^{\cW_4}= (13/6)\cdot 2=13/3\), the required next coefficient is
\[
 A_5^{\cW_4}
 =
 -\frac{52}{5}\cdot \frac{13}{3}
 =
 -\frac{676}{15}.
\]
Thus the bridge is not proved by the present lane computation: the
lane computation supplies \(13\), while the full Miura/OPE computation
of \(A_5^{\cW_4}\) is absent.
\end{remark}

\section{The closed form $\beta_N = 12(H_N - 1)$}
\label{sec:beta-N-closed-form-theorem}

\begin{conjecture}[Harmonic closed form of $\beta_N$]
\label{conj:beta-N-harmonic-closed-form}
\ClaimStatusConjectured
For every integer $N \ge 2$, the Riccati leading-Laurent ratio coefficient
$\beta_N$ of the principal $\cW_N$ shadow tower satisfies
\begin{equation}
\label{eq:beta-N-closed-form}
\beta_N
\;=\;
12 \cdot (H_N - 1)
\;=\;
12 \sum_{j = 2}^{N} \frac{1}{j},
\end{equation}
where $H_N = \sum_{j=1}^{N} 1/j$ is the $N$-th harmonic number. Explicit
values:
\[
 \beta_2 = 6,\quad \beta_3 = 10,\quad \beta_4 = 13,\quad
 \beta_5 = \tfrac{77}{5},\quad \beta_6 = \tfrac{87}{5},\quad
 \beta_7 = \tfrac{669}{35}, \ldots
\]
\end{conjecture}

\begin{remark}[Conditional derivation from the scaling law]
By Conjecture~\ref{conj:kappa-ratio-scaling-law},
\[
 A_r(\cW_N) \;=\; \bigl(2(H_N - 1)\bigr)^{r-3} \cdot A_r(\mathrm{Vir}),
 \qquad A_r(\mathrm{Vir}) = \frac{8 (-6)^{r-4}}{r}.
\]
Taking the ratio at $r+1$ and $r$:
\begin{align*}
 \frac{A_{r+1}(\cW_N)}{A_r(\cW_N)}
 &\;=\;
 \bigl(2(H_N - 1)\bigr)^{(r+1-3) - (r-3)}
 \cdot \frac{A_{r+1}(\mathrm{Vir})}{A_r(\mathrm{Vir})} \\
 &\;=\;
 2(H_N - 1) \cdot
 \Bigl( -\frac{6 r}{r + 1} \Bigr) \\
 &\;=\;
 -\frac{12 (H_N - 1) \cdot r}{r + 1}.
\end{align*}
Comparing with the defining form $A_{r+1}/A_r = -\beta_N \cdot r / (r + 1)$
(Lemma~\ref{lem:leading-coefficient-ratio-identity} at general $N$):
\[
 \beta_N \;=\; 12 (H_N - 1).
\]

Verification at the proved data:
\begin{itemize}[nosep]
\item $N = 2$: $H_2 - 1 = 1/2$, so $\beta_2 = 12 \cdot 1/2 = 6$,
 matching Lemma~\ref{lem:leading-coefficient-ratio-identity}.
\item $N = 3$: $H_3 - 1 = 5/6$, so $\beta_3 = 12 \cdot 5/6 = 10$,
 matching Proposition~\ref{prop:w3-shadow-leading-asymptotic}(4).
\end{itemize}
\end{remark}

\begin{corollary}[Polynomial extrapolations excluded]
\label{cor:candidate-A-B-retracted}
\ClaimStatusConditional
Both the triangular formula $\beta_N^{(A)} = (N+1)(N+2)/2$ and the
quadratic formula $\beta_N^{(B)} = N^2 - N + 4$ of Chapter~\ref{ch:wn-tempered-closure}
(Proposition~\ref{prop:beta-N-closed-form} and
Remark~\ref{rem:beta-candidate-discriminator})
are conditionally excluded. At $N = 4$ the harmonic value is $\beta_4 = 13$, whereas
the triangular formula gives $15$ and the quadratic formula gives $16$.
The agreement at $N = 2, 3$ between the two polynomial formulas and
the harmonic candidate $\beta_N = 12(H_N-1)$
is a \emph{coincidence}: both candidates are quadratic in $N$, and the
harmonic formula $12(H_N - 1)$ happens to equal a quadratic in $N$ at exactly
two data points before diverging.
\end{corollary}

\begin{proof}
Conditional on Conjecture~\ref{conj:beta-N-harmonic-closed-form},
direct evaluation at $N = 4$ gives:
\begin{align*}
 \beta_4^{(A)} &\;=\; \tfrac{5 \cdot 6}{2} = 15, \\
 \beta_4^{(B)} &\;=\; 16 - 4 + 4 = 16, \\
 \beta_4^{\mathrm{harm}} &\;=\; 12 \cdot (H_4 - 1)
 = 12 \cdot \tfrac{13}{12}
 = 13.
\end{align*}
All three values are pairwise distinct, and the harmonic conjecture
(Conjecture~\ref{conj:beta-N-harmonic-closed-form}) selects
the last. The agreement at $N = 2, 3$ is interpolation without force:
two values do not determine a quadratic polynomial. The triangular and
quadratic formulas are two different polynomial continuations through
the same two low-rank points, and both miss the harmonic value at the
first unfixed rank.
\end{proof}

\begin{corollary}[$\beta_N$ rational and non-integral for \(N\geq 5\)]
\label{cor:beta-N-rational-not-integer}
\ClaimStatusConditional
Conditional on the closed form $\beta_N = 12(H_N - 1)$, one has
\[
 \beta_2 = 6,\qquad \beta_3 = 10,\qquad \beta_4 = 13
\]
in $\mathbb Z$, and $\beta_N \notin \mathbb Z$ for every $N\geq 5$.
The first non-integral value is $\beta_5 = 77/5$.  If
\[
 h_N := \operatorname{den}(H_N-1),\qquad
 d_N := \operatorname{den}(\beta_N),
\]
then
\[
 d_N = \frac{h_N}{\gcd(h_N,12)}.
\]
In particular $d_N$ divides
$\mathrm{lcm}(1,2,\ldots,N)/\gcd(\mathrm{lcm}(1,2,\ldots,N),12)$, but
the exact-denominator equality
$d_N=\mathrm{lcm}(1,2,\ldots,N)/12$ is false in general.  It already
fails at $N=18$:
\[
 H_{18}-1=\frac{10190221}{4084080},\qquad
 \beta_{18}=\frac{10190221}{340340},\qquad
 \frac{\mathrm{lcm}(1,\ldots,18)}{12}=1021020.
\]
\end{corollary}

\begin{proof}
Conditional on Conjecture~\ref{conj:beta-N-harmonic-closed-form},
the rationality and the values $6,10,13,77/5$ are direct evaluations.
Write $H_N-1=a_N/h_N$ in lowest terms.  Then
\[
 \beta_N=12a_N/h_N,
\]
so reduction gives
$d_N=h_N/\gcd(h_N,12)$.  Since
$h_N\mid \mathrm{lcm}(1,\ldots,N)$, this gives the displayed
divisibility.

It remains to prove non-integrality for all $N\geq 5$.  By Bertrand's
postulate, applied to $m=\lceil N/2\rceil$, there is a prime $p$ with
$N/2<p\leq N$; for $N\geq 5$ this prime may be chosen with $p\geq 5$.
Let $L_N=\mathrm{lcm}(1,\ldots,N)$ and set
\[
 A_N=L_N\sum_{j=2}^{N}\frac1j.
\]
Because $p>N/2$, the only multiple of $p$ in $\{1,\ldots,N\}$ is $p$
itself, hence $v_p(L_N)=1$.  In the sum defining $A_N$, the $j=p$
term contributes $L_N/p$, which is not divisible by $p$, while every
term with $j\neq p$ contributes a multiple of $p$.  Thus
$A_N\not\equiv 0\pmod p$, so $p$ survives in the reduced denominator
of $H_N-1$.  Since $p\nmid 12$, it also survives in the reduced
denominator of $\beta_N=12(H_N-1)$.  Therefore $\beta_N$ is not an
integer for every $N\geq 5$.

The numerical assertion at $N=18$ follows by exact arithmetic:
$\mathrm{lcm}(1,\ldots,18)=12252240$ and
$12(10190221/4084080)=10190221/340340$.
\end{proof}

\begin{proposition}[Convergence radius closed form;
\ClaimStatusConditional; licensing $\alpha+\gamma$ via fixed rank and
central charge, Conjecture~\ref{conj:beta-N-harmonic-closed-form},
leading-root sharpness, and $\hypAmbientWtCpl$]
\label{prop:rho-star-WN-closed-form-upgrade}
\ClaimStatusConditional
Under Conjecture~\ref{conj:beta-N-harmonic-closed-form}, assume
leading-root sharpness of the ordinary shadow sequence,
\[
 \limsup_{r \to \infty} \lvert S_r(\cW_{N,c}) \rvert^{1/r}
 =
 \beta_N/\lvert c\rvert .
\]
For every $N \ge 2$ and $c \in \bbR \setminus (\{0\} \cup
\cD^{\mathrm{sev}}_{\mathrm{Kac}}(\cW_N))$ satisfying this sharpness
hypothesis,
\[
 \rho_*^{\cW_N}(c)
 \;=\;
 \frac{\lvert c \rvert}{\beta_N}
 \;=\;
 \frac{\lvert c \rvert}{12 (H_N - 1)}
 \;=\;
 \frac{\lvert c \rvert}{12 \sum_{j = 2}^{N} 1/j}.
\]
Specialisations:
$\rho_*^{\mathrm{Vir}}(c) = |c|/6$ ($N=2$, matching
Proposition~\ref{prop:virasoro-rho-star-closed-form});
$\rho_*^{\cW_3}(c) = |c|/10$ ($N=3$, matching
Theorem~\ref{thm:tempered-stratum-contains-w3});
$\rho_*^{\cW_4}(c) = |c|/13$;
$\rho_*^{\cW_5}(c) = 5|c|/77$.
\end{proposition}

\begin{proof}
The Cauchy--Hadamard root test gives
\[
 \rho_*^{\cW_N}(c)
 =
 \Bigl(\limsup_{r\to\infty}
 \lvert S_r(\cW_{N,c})\rvert^{1/r}\Bigr)^{-1}.
\]
The sharpness hypothesis identifies the limsup with $\beta_N/|c|$.
Substituting \eqref{eq:beta-N-closed-form} gives the displayed formula.
Without sharpness, the finite envelope gives only
$\rho_*^{\cW_N}(c)\ge |c|/\beta_N$.
\end{proof}

\section{Numerical values through $N = 12$}
\label{sec:beta-N-extended-table}

\begin{table}[h]
\centering
\caption{Closed form $\beta_N = 12(H_N - 1)$ across $N = 2, \ldots, 12$.
Columns: $H_N$, $\beta_N$ (exact rational), $\beta_N$ (decimal), and
the triangular/quadratic polynomial interpolants.}
\label{tab:beta-N-values-extended}
\renewcommand{\arraystretch}{1.2}
\begin{tabular}{|c|c|l|l|c|c|}
\hline
$N$ & $H_N$ & $\beta_N$ (exact) & $\beta_N$ (dec.) & Tri. & Quad. \\
\hline
$2$ & $3/2$ & $6$ & $6.000$ & $6$ & $6$ \\
$3$ & $11/6$ & $10$ & $10.000$ & $10$ & $10$ \\
$4$ & $25/12$ & $13$ & $13.000$ & $15$ & $16$ \\
$5$ & $137/60$ & $77/5$ & $15.400$ & $21$ & $24$ \\
$6$ & $49/20$ & $87/5$ & $17.400$ & $28$ & $34$ \\
$7$ & $363/140$ & $669/35$ & $19.114$ & $36$ & $46$ \\
$8$ & $761/280$ & $1443/70$ & $20.614$ & $45$ & $60$ \\
$9$ & $7129/2520$ & $4609/210$ & $21.948$ & $55$ & $76$ \\
$10$ & $7381/2520$ & $4861/210$ & $23.148$ & $66$ & $94$ \\
$11$ & $83711/27720$ & $55991/2310$ & $24.239$ & $78$ & $114$ \\
$12$ & $86021/27720$ & $58301/2310$ & $25.239$ & $91$ & $136$ \\
\hline
\end{tabular}
\end{table}

\begin{remark}[Asymptotic $\beta_N \sim 12 \log N$]
\label{rem:beta-asymptotic}
From $H_N \sim \log N + \gamma$ (Euler-Mascheroni constant
$\gamma \approx 0.5772$), one has
\[
 \beta_N
 \;=\;
 12(H_N - 1)
 \;\sim\;
 12 (\log N + \gamma - 1)
 \;\approx\;
 12 \log N - 5.07.
\]
This logarithmic growth in $N$ contrasts sharply with the quadratic
growth predicted by the triangular formula ($\beta_N \sim N^2/2$) and
the quadratic formula ($\beta_N \sim N^2$). At $N = 100$,
$\beta_{100} \approx 12 \cdot 4.61 \approx 55.3$, whereas the
triangular formula gives $\beta_{100}^{(A)} =
5151$ (a factor $\sim 93$ larger). If the harmonic law holds,
$\beta_N$ grows much more slowly than any polynomial guess.
\end{remark}

\section{Tempering theorem with conditional exact radius}
\label{sec:tempering-theorem-healed}

\begin{theorem}[$\cW_N$ tempering with conditional exact radius;
\ClaimStatusConditional; licensing $\alpha+\gamma+\varepsilon$ via fixed
rank, finite Riccati envelope, leading-root sharpness for the radius,
finite propagation, and $\hypAmbientWtCpl$]
\label{thm:wn-tempered-exact-radius}
\ClaimStatusConditional
For every integer $N \ge 2$ and every $c \in \bbR \setminus (\{0\} \cup
\cD^{\mathrm{sev}}_{\mathrm{Kac}}(\cW_N))$, assume the principal
$\cW_N$ shadow tower has a finite Riccati envelope. Then the principal
$\cW_N$ algebra $\cW_{N,c}$ satisfies the analytic tempering hypothesis
$(\mathrm{H}_S)$:
\[
 \limsup_{r \to \infty}
 \bigl( \lvert S_r(\cW_{N,c}) \rvert / r! \bigr)^{1/r}
 \;=\;
 0,
\]
and the chain-level $E_3$-topological structure on the original
curved-Dunn complex
(Theorem~\ref{thm:chain-level-e3-original-complex-conditional}) has
its analytic tempering hypothesis discharged for $\cW_{N,c}$ on this
finite-envelope locus. If Conjecture~\ref{conj:beta-N-harmonic-closed-form}
holds and the finite envelope is leading-root sharp, the
ordinary-generating convergence radius has the closed form
\[
 \rho_*^{\cW_N}(c)
 \;=\;
 \frac{|c|}{12 \sum_{j=2}^{N} 1/j}
 \;=\;
 \frac{|c|}{12(H_N - 1)}.
\]
\end{theorem}

\begin{proof}
Combine Theorem~\ref{thm:wn-tempered-all-N} (finite envelope plus
Stirling dominance) with Conjecture~\ref{conj:beta-N-harmonic-closed-form}
(exact closed form $\beta_N = 12(H_N - 1)$). The Stirling argument is
$\beta$-independent, so the tempering conclusion follows from the finite
envelope before the exact arithmetic of $\beta_N$ is used. For the
ordinary radius, Proposition~\ref{prop:rho-star-WN-closed-form-upgrade}
adds leading-root sharpness; then substitution gives
$|c|/(12(H_N - 1))$.
\end{proof}

\section{Arithmetic evidence}
\label{sec:beta-N-independent-verification}

Three scaling-ansatz evidence paths and a separate finite-rank
$\cW_4$ calculation test the harmonic scaling hypothesis rather than
prove it:

\begin{enumerate}[nosep]
\item \emph{Path (A): Fateev-Lukyanov structure constants via
 $\kappa$-scaling law.}
 Derived from: Vol II datum $A_4^{\cW_3} = 10/3$ plus
 $\kappaChHodge(\cW_N) = c(H_N - 1)$.
 Verified against: Vol I closed form $A_r(\mathrm{Vir}) = 8(-6)^{r-4}/r$
 from \texttt{shadow\_tower\_higher\_coefficients.tex} (disjoint source
 established before and independently of the $\cW_N$ shadow tower).

\item \emph{Path (B): Laurent arithmetic under the kappa-scaling ansatz.}
 Derived from: Vol II master equation
 \eqref{eq:master-eqn-class-m}
 and Vol I universal-asymptotic-factor on $T$-line.
 Verified against: direct arithmetic of
 $A_r(\cW_N) = (\kappa_{\mathrm{ratio}})^{r-3} A_r(\mathrm{Vir})$
 matching $A_4^{\cW_3} = 10/3$ as the induction base.

\item \emph{Path (C): Riccati-analog recurrence under the harmonic ansatz.}
 Derived from: conditional formulas for $\beta_N$ and $A_r$ from
 Conjecture~\ref{conj:beta-N-harmonic-closed-form}.
 Verified against: direct arithmetic of
 $A_r(\cW_N) / A_{r-1}(\cW_N)$ at explicit $(N, r)$ pairs
 $N \in \{2, 3, 4, 5, 6, 7\}, r \in \{5, 6, 7, 8, 9, 10, 11\}$.
 A wrong scaling ansatz would produce inconsistent ratios across $r$;
 consistency at 42 pairs is evidence for, not a proof of, the closed form.

\item \emph{Path (D): Direct $\cW_4$ spin-lane calculation.}
 Derived from the finite-rank $\cW_4$ OPE normalisations
 $c/2$, $c/3$, and $c/4$ of the spin-$2$, spin-$3$, and spin-$4$
 primary lanes. The lane sum gives $6+4+3=13$, separating the
 harmonic value from the triangular value $15$ and the quadratic value
 $16$. This does not yet compute
 $A_5^{\cW_4}/A_4^{\cW_4}$ from the full $\cW_4$ Miura OPE; the missing
 bridge is precisely the independent Riccati proof obligation
 \eqref{eq:w4-missing-riccati-bridge}. Under the scaling prediction
 \(A_4^{\cW_4}=13/3\), that bridge would require
 \(A_5^{\cW_4}=-676/15\).
\end{enumerate}

\emph{Test coverage.} The beta and $\cW_4$ tests cover:
\begin{itemize}[nosep]
\item Basic unit tests for $H_N$, $\kappa$-ratio (3 tests).
\item Path (A) tests: known data, $\cW_2 = \mathrm{Vir}$ collapse,
 $\beta_4$ prediction (3 tests, 1 independent-verification decorated).
\item Path (B) tests: $A_4^{\cW_3}$, $A_5^{\cW_3}$, and the conditional
 $A_5^{\cW_4}$ value under the kappa-scaling ansatz
 (3 tests, 1 independent-verification decorated).
\item Path (C) tests: ratio identity at all $(N, r)$, explicit $N=4$
 match (2 tests, 1 independent-verification decorated).
\item Discrimination tests: under the harmonic ansatz, exclude the
 triangular and quadratic formulas at $N=4$ and detect the first
 rational non-integer value at $N=5$; exact arithmetic guards reject
 the false lcm-denominator formula and verify the Bertrand-prime
 denominator obstruction for $N\geq 5$ (7 tests).
\item Extrapolation sanity: monotonicity, $\log N$ asymptotic (2 tests).
\item Direct $\cW_4$ tests: spin-lane norms, lane sum
 $13$, the required Riccati ratio \(-52/5\), and the explicit record
 that the full-Miura \(A_5^{\cW_4}\) computation remains open.
\end{itemize}

\section{Polynomial extrapolations excluded}
\label{sec:excluded-polynomial-extrapolations}

The following polynomial extrapolations agree with the low-rank data
only accidentally and are incompatible with the harmonic candidate
conditional on Conjecture~\ref{conj:beta-N-harmonic-closed-form}:

\begin{enumerate}[nosep]
\item The triangular formula $(N+1)(N+2)/2$: conditionally excluded at
 $N \ge 4$. The coincidences at $N = 2, 3$ remain correct because
 they are also consequences of
 Conjecture~\ref{conj:beta-N-harmonic-closed-form}.

\item The quadratic formula $N^2 - N + 4$: conditionally excluded at
 $N = 4$. Both candidates were coincidences of low-rank interpolation
 through $N = 2, 3$.

\item Any sharp convergence-radius formula using a polynomial denominator:
 the harmonic candidate gives $|c|/[12(H_N - 1)]$ after leading-root
 sharpness is imposed.
\end{enumerate}

Theorem~\ref{thm:wn-tempered-all-N} of Chapter~\ref{ch:wn-tempered-closure}
depends only on the finite Riccati envelope and $\beta$-independent
Stirling dominance, not on the exact closed form of $\beta_N$.

\section{Harmonic consequences}
\label{sec:beta-N-platonic-consequences}

The harmonic candidate $\beta_N = 12(H_N - 1)$ has three
conditional consequences.

\emph{(1) Harmonic numbers are the universal $\cW_N$ invariant.}
The formula $\kappaChHodge(\cW_N) = c(H_N - 1)$ (Vol~I kappa formula) and the
formula $\beta_N = 12(H_N - 1)$ (this conjecture) differ only in the
scalar prefactor. Both invariants are controlled by the same harmonic
sum $H_N - 1 = \sum_{j=2}^{N} 1/j$, which is the principal $\cW_N$
line-decomposition data: the sum $1/j$ counts the curvature $\kappa_\ell =
c/j$ of the $\cW^{(j)}$-line for $j = 2, \ldots, N$. The $\cW_N$
shadow tower's sharp analytic radius is the harmonic average of its
line curvatures, conditional on the scaling law and leading-root
sharpness.

\emph{(2) The coincidence at $N = 2, 3$ is quadratic-interpolation
generic.} Any polynomial-looking closed form that matches two data
points has multiple candidates; the coincidence with $(N+1)(N+2)/2$
and $N^2 - N + 4$ at $N = 2, 3$ is a consequence of both being
quadratic in $N$ with appropriate offsets. The harmonic formula
$12(H_N - 1)$ is not polynomial in $N$, so every polynomial candidate
eventually diverges.

\emph{(3) The logarithmic $N$-growth is physically natural if the
harmonic law holds.}
$\beta_N \sim 12 \log N$ at large $N$ reflects the fact that the sharp
radius scale $|c|/\beta_N$ for higher-rank $\cW_N$ algebras slowly
shrinks as $N$ increases. This is consistent with the
$\cW_\infty[\mu]$ inverse-limit construction:
$\beta_\infty = \lim_{N \to \infty} 12 (H_N - 1) = \infty$, so this
sharp-radius scale tends to zero at every fixed $c$ and the finite-rank
argument supplies no uniform $N$-independent Banach radius. The
shadow-tower analytic structure degenerates at $\cW_\infty[\mu]$,
consistent with
Remark~\ref{V1-rem:gaberdiel-gopakumar-truncation}
(Vol~I \texttt{shadow\_tower\_extended\_families.tex}) on the
Gaberdiel-Gopakumar truncation loci.

\begin{remark}[anchor: $\beta_N$ as the rung-$\cW_N$ /
rung-$\cW_\infty$ tempering rate]%
\label{rem:anchor-beta-N-rung-w-N-w-infty}%
%
%
\index{$\beta_N$!chain-level avatar at every rung}%
Conditional on Conjecture~\ref{conj:beta-N-harmonic-closed-form},
the closed form $\beta_N = 12(H_N - 1)$ is the precise
\emph{tempering rate} attached to the rung-$\cW_N$ corollary
(\textsf{cor:rung-w-N}) of the universal-holography master theorem
(\textsf{thm:universal-holography-master} in
\textsf{chap:universal-conductor}; \textsf{chap:climax-platonic}). The
linkage:
\begin{itemize}[nosep]
\item Per \textsf{cor:rung-w-N}, the iterated Sugawara construction
 (Vol~II \texttt{e\_infinity\_topologization.tex},
 \textsf{thm:iterated-sugawara-construction}) lifts the chiral
 Higher-Deligne $E_3$-action on $Z^{\mathrm{der}}_{\mathrm{ch}}$ to
 $E_{N+1}$-topological for principal $\cW_N$. The chain-level
 closure of this lift requires the tempering hypothesis
 $(\mathrm{H}_S)$ on the curved-Dunn complex
 (Theorem~\ref{thm:chain-level-e3-original-complex-conditional}); the
 hypothesis discharges by finite-envelope Stirling dominance; the sharp
 convergence radius, when leading-root sharpness is imposed, is
 $\rho_*^{\cW_N}(c) = |c|/\beta_N = |c|/[12(H_N - 1)]$
 (Proposition~\ref{prop:rho-star-WN-closed-form-upgrade}).
\item Per \textsf{cor:rung-w-infty}, the inverse limit
 $\cW_\infty = \varprojlim_N \cW_N$ realizes the $E_\infty$-topological
 horizon. The vanishing of the sharp-radius scale
 ($\beta_\infty = \infty$, item~(3) above) means the $\cW_\infty$
 chain-level rung is tempered only in the per-truncation
 $\rho_*^{\cW_N}$ sense, not at uniform $\rho$ as $N \to \infty$. The
 horizon is reached as the limit of finite truncations, not as a
 single algebra at infinite rank with uniform Banach norm. The
 $E_\infty$-topological climax is therefore a \emph{tower limit}, not
 a chain-level identity at $\cW_\infty[\mu]$.
\item Conditional on the harmonic beta law, the identity
 $\kappaChHodge(\cW_N) = c(H_N - 1)$ shared by the chiral central curvature
 and the tempering rate
 ($\beta_N = 12(H_N - 1)$) collects the climactic
 $\cW_N$-line geometry into a single arithmetic invariant: the
 $H_N - 1$ harmonic sum is the \emph{universal $\cW_N$ conductor}
 of the climax (cf.\ \textsf{chap:universal-conductor}).
\end{itemize}
This remark records the structural anchor: conditional on the scaling
law, the harmonic closed form sets the finite-rung sharp convergence
radius after leading-root sharpness is imposed, while the infinite-rank
horizon is a tower limit.
\end{remark}

\section{Summary}
\label{sec:beta-N-summary}

\begin{itemize}[nosep]
\item Harmonic closed form conjectured for every $N \ge 2$:
 $\beta_N = 12(H_N - 1) = 12 \sum_{j=2}^{N} 1/j$
 (Conjecture~\ref{conj:beta-N-harmonic-closed-form}).
\item Conditional value $\beta_4 = 13$; the triangular and quadratic interpolants give
 $15$ and $16$
 (Corollary~\ref{cor:candidate-A-B-retracted}).
\item Conditional on the harmonic law, $\beta_N$ is rational, not integer, for $N \ge 5$ ($\beta_5 = 77/5$,
 first non-integer).
\item Conditional sharp convergence radius:
 $\rho_*^{\cW_N}(c) = |c|/[12(H_N - 1)]$
 after leading-root sharpness is imposed
 (Proposition~\ref{prop:rho-star-WN-closed-form-upgrade}).
\item Tempering theorem: every finite $\beta_N$ is absorbed by
 Stirling dominance.
\item The three scaling-ansatz paths are disjoint from the direct
 $\cW_4$ finite-rank calculation.
\end{itemize}

For every finite principal $\cW_N$, the tempering statement
(Chapter~\ref{ch:wn-tempered-closure}) is proved at every finite
$\beta_N$; the exact harmonic convergence radius in this chapter is
conditional on the $\kappa$-ratio scaling law and leading-root
sharpness.
